% vision_standalone.tex — Comprehensive standalone vision system report
% SAR Drone Project, University of Bristol AENGM0074
% Compile standalone or \input{} into main report

\section{Vision System Architecture}
\label{sec:vision-arch}

The vision subsystem is the sole module responsible for camera capture, AI inference,
and object detection across the entire SAR drone codebase. All detection logic is
encapsulated within a single file (\texttt{vision.py}), exposing a unified interface
that every other module consumes without knowledge of the underlying backend.

\subsection{Design Philosophy}

A strict separation-of-concerns principle governs the vision architecture: only
\texttt{vision.py} touches computer vision and AI libraries. Every other module---the
state machine (\texttt{main.py}), the ground station (\texttt{pi\_flight.py}), the
passive observer (\texttt{passive\_watch.py}), and all test scripts---calls a single
function:

\begin{verbatim}
found, cx, cy, conf = eyes.detect_in_image(frame)
\end{verbatim}

\noindent where \texttt{found} is a boolean, \texttt{cx} and \texttt{cy} are pixel
coordinates of the detection centre, and \texttt{conf} is a float in $[0, 1]$.
This interface contract means that any change to the model, preprocessing pipeline,
inference backend, or postprocessing logic requires zero modifications to calling code.

\subsection{Dual-Backend Architecture}

The \texttt{VisionSystem} class supports three inference backends, selected
automatically at initialisation based on available libraries and user preference:

\begin{enumerate}
    \item \textbf{Ultralytics YOLO} --- Full-featured backend for development machines
          with PyTorch and CUDA. Supports \texttt{.pt} model files directly. Used on
          the Windows laptop during simulation and video analysis.
    \item \textbf{TFLite} --- Lightweight backend for the Raspberry Pi~5.
          Supports \texttt{.tflite} model files via \texttt{tflite-runtime},
          \texttt{ai-edge-litert} (Python~3.13 compatible), or TensorFlow Lite.
          The primary deployment backend.
    \item \textbf{NCNN} --- ARM-optimised backend from Tencent, providing approximately
          $3\times$ speedup over TFLite on the Pi~5 (13.8~FPS vs 4.8~FPS). Selectable
          via the \texttt{-{}-backend ncnn} CLI flag or \texttt{backend="ncnn"} parameter.
\end{enumerate}

Backend selection follows a priority chain: NCNN (if explicitly requested) $\rightarrow$
Ultralytics (if available and model is \texttt{.pt}) $\rightarrow$ TFLite (fallback).
This auto-detection is performed once at import time by probing for each library:

\begin{verbatim}
try:
    from ultralytics import YOLO as _YOLO
except ImportError:
    pass
try:
    from tflite_runtime.interpreter import Interpreter
except ImportError:
    try:
        from ai_edge_litert.interpreter import Interpreter
    except ImportError:
        pass
\end{verbatim}

Adding a new backend requires implementing two methods---\texttt{\_try\_load\_*} and
\texttt{\_detect\_*}---with zero changes to any caller. This extensibility was validated
when NCNN was added as the third backend six weeks after the initial dual-backend design.

\subsection{Threaded Inference Pipeline}

On the Raspberry Pi, TFLite inference takes approximately 206~ms per frame. A
single-threaded approach would limit both the video stream and detection rate to
4.8~FPS, producing an unacceptable stuttering stream for the ground station operator.

The solution decouples capture from inference using a background thread
(Design Decision DD-11):

\begin{itemize}
    \item \textbf{Main thread (display loop):} Grabs camera frames at native rate
          ($\sim$30~FPS on Pi), draws overlays (GPS crosshair, detection boxes from
          the last result), encodes JPEG, and serves the MJPEG stream.
    \item \textbf{Inference thread:} Grabs the latest frame from a shared variable,
          runs \texttt{detect\_in\_image()}, and updates shared detection state.
          Frame-dropping is intentional: inference always processes the most recent
          frame, skipping stale ones.
\end{itemize}

Shared state is protected by two locks: \texttt{\_inference\_lock} for detection
results and \texttt{\_inference\_eyes\_lock} for safe model hot-swapping. The
detection overlay lags by one inference cycle ($\sim$200~ms), which is imperceptible
at flight altitude.


%% ======================================================================
\section{Camera and Sensor}
\label{sec:camera-sensor}

\subsection{IMX296 Global Shutter Sensor}

The drone carries a Sony IMX296 global shutter camera interfaced via the Raspberry
Pi Camera Module~3 (or compatible). Key specifications are summarised in
Table~\ref{tab:sensor-specs}.

\begin{table}[htbp]
\centering
\caption{IMX296 sensor and camera configuration.}
\label{tab:sensor-specs}
\begin{tabular}{ll}
\hline
\textbf{Parameter}       & \textbf{Value} \\
\hline
Sensor                   & Sony IMX296 \\
Shutter type             & Global shutter \\
Native resolution        & 1456 $\times$ 1088 pixels \\
Sensor width             & 5.02~mm \\
Focal length (calibrated)& 5.46~mm \\
Horizontal FOV           & $\sim$49.3\textdegree \\
Vertical FOV             & $\sim$36.8\textdegree \\
Interface                & CSI-2 via \texttt{picamera2} \\
Colour format            & BGR (despite RGB888 label) \\
Camera flip              & 180\textdegree\ (mounted inverted) \\
\hline
\end{tabular}
\end{table}

\subsection{BGR Colour Discovery}

A persistent blue colour cast was observed in early testing despite adjusting
auto white balance modes, manual colour gains, and ISP tuning parameters.
The root cause was identified by testing all six permutations of channel ordering:
the IMX296 sensor outputs BGR data despite \texttt{picamera2} labelling the
format as RGB888. The codebase initially included a
\texttt{cv2.cvtColor(frame, cv2.COLOR\_RGB2BGR)} conversion that double-swapped
the channels.

The fix was to remove the \texttt{cvtColor} call entirely. This discovery
(Design Decision DD-06) underscores that sensor format labels should be
validated empirically rather than trusted.

\subsection{Camera Initialisation}

\texttt{VisionSystem} attempts OpenCV \texttt{VideoCapture} first (for laptop
webcams), then falls back to \texttt{picamera2} for the Pi camera. On Windows,
the DirectShow backend (\texttt{cv2.CAP\_DSHOW}) is used; on Linux, the default
V4L2 backend is selected. The camera is configured at the native 1456$\times$1088
resolution with a target of 30~FPS. Ten warm-up frames are discarded to allow AWB
and auto-exposure convergence.


%% ======================================================================
\section{AI Model}
\label{sec:ai-model}

\subsection{Architecture: YOLOv8n}

All deployed models use the YOLOv8 nano architecture~\cite{yolov8},
the smallest variant in the YOLOv8 family. It provides a single-class detector
for the ``dummy'' target with the following tensor shapes:

\begin{table}[htbp]
\centering
\caption{YOLOv8n model specifications.}
\label{tab:model-specs}
\begin{tabular}{lll}
\hline
\textbf{Property} & \textbf{TFLite} & \textbf{NCNN} \\
\hline
Input shape       & $[1, 640, 640, 3]$ (NHWC) & $[1, 3, 640, 640]$ (NCHW) \\
Input dtype       & float32, $[0, 1]$ normalised & float32 \\
Output shape      & $[1, 5, 8400]$ & $[1, 5, 8400]$ \\
Output format     & $(x, y, w, h, \text{conf})$ per anchor & same \\
File size          & $\sim$11.7~MB & $\sim$7~MB \\
Classes           & 1 (dummy) & 1 (dummy) \\
\hline
\end{tabular}
\end{table}

The output tensor contains 8400 anchor predictions, each with five values:
bounding box centre $(x, y)$, width $w$, height $h$, and confidence score.
Non-maximum suppression (NMS) is applied in \texttt{vision.py} to select the
highest-confidence detection above the threshold.

\subsection{Model Evolution}

The model was iteratively improved across three training generations, summarised
in Table~\ref{tab:model-history}.

\begin{table}[htbp]
\centering
\caption{Model training history.}
\label{tab:model-history}
\begin{tabular}{llllll}
\hline
\textbf{Version} & \textbf{Dataset} & \textbf{imgsz} & \textbf{mAP50} & \textbf{Size} & \textbf{Notes} \\
\hline
v1 (original) & 200 synthetic, 640$\times$640 & 640 & $\sim$0.95 & 3.3~MB & Single background \\
sar\_640       & Improved synthetic & 640 & --- & --- & Experiment \\
sar\_1280      & Improved synthetic & 1280 & --- & --- & Experiment \\
sar\_v2\_1088  & 300 syn + 16 real + 50 neg & 1088 & 0.995 & 11.7~MB & Current best \\
sar\_v3 (planned) & 500 syn + 50+ real + 100 neg & 640/1088 & TBD & --- & Next iteration \\
\hline
\end{tabular}
\end{table}

\textbf{v1} used 200 synthetic composite images generated by pasting
\texttt{dummy.png} onto a single satellite image (\texttt{map.jpg}) at
640$\times$640. While sufficient to prove the detection pipeline, the model
memorised the background texture and struggled with real outdoor scenes.

\textbf{v2 (sar\_v2\_1088)} addressed three critical v1 limitations:
\begin{itemize}
    \item \textbf{Multiple backgrounds:} 300 synthetic images composited onto
          random crops from DJI flight video frames, preventing background overfitting.
    \item \textbf{Real labelled frames:} 16 frames from actual flight footage,
          labelled at native 1456$\times$1088 using \texttt{tools/label\_tool.py -{}-full}.
    \item \textbf{Negative images:} 50 frames with no target (empty label files),
          teaching the model what ``no dummy'' looks like and reducing false positives.
\end{itemize}

Training was performed on Google Colab with a T4 GPU: \texttt{imgsz=1088},
\texttt{epochs=150}, \texttt{batch=8}, with early stopping at \texttt{patience=30}.
The resulting model achieved mAP50 = 0.995, a substantial improvement over v1.

\subsection{Training Resolution vs Inference Resolution}

An important design choice is the mismatch between training resolution (1088) and
inference resolution (640). Training at the native camera resolution allows the
model to learn features at the actual pixel scale present in real frames. However,
the TFLite export is fixed at 640$\times$640 input because \texttt{vision.py}
resizes every frame to this size before inference. This is a well-established
pattern in the YOLO ecosystem: training at higher resolution improves feature
learning while the runtime handles downscaling transparently.


%% ======================================================================
\section{Data Augmentation and Synthetic Generation}
\label{sec:data-augmentation}

\subsection{Synthetic Dataset Generation Pipeline}

The script \texttt{generate\_dataset\_v2.py} creates training data by compositing
a dummy foreground image (\texttt{dummy.png} with alpha channel) onto background
crops from real DJI flight video frames. The generation process:

\begin{enumerate}
    \item Extracts random background crops from DJI video at 1456$\times$1088.
    \item Scales the dummy image to altitude-correct pixel size using:
          \[
          \text{apparent\_size\_px} = \frac{\text{real\_height\_m} \times f_{\text{px}}}
          {\text{altitude\_m}}
          \]
          At 30~m altitude with $f_{\text{px}} = 1416$ and dummy height 1.8~m, this gives
          approximately 85 pixels.
    \item Applies Gaussian blur to the alpha mask edges (7$\times$7 kernel) to prevent
          sharp composite boundaries that the model could exploit.
    \item Composites the dummy at a random position within the background.
    \item Generates YOLO-format label files (class, $x_\text{centre}$, $y_\text{centre}$,
          width, height, all normalised to $[0, 1]$).
\end{enumerate}

\subsection{Augmentation Types}

Table~\ref{tab:augmentations} ranks the augmentation techniques by impact for
aerial SAR detection.

\begin{table}[htbp]
\centering
\caption{Data augmentation techniques ranked by impact for aerial SAR.}
\label{tab:augmentations}
\begin{tabular}{clp{6.5cm}}
\hline
\textbf{Rank} & \textbf{Augmentation} & \textbf{Rationale} \\
\hline
1 & Motion blur          & Drone movement at 5~m/s causes directional blur; the primary
                           failure mode in real flights. Kernel size 5--15~px. \\
2 & Brightness/contrast  & Lighting changes drastically between conditions.
                           HSV jitter: $h{=}0.015$, $s{=}0.7$, $v{=}0.4$. \\
3 & Gaussian noise        & Sensor noise at high shutter speeds. $\sigma = 5$--$20$. \\
4 & Shadow overlays       & Trees, buildings, and the drone itself cast shadows.
                           Dark rectangles/circles composited before dummy. \\
5 & Perspective transforms& 5--15\textdegree\ tilt during flight and wind gusts.
                           \texttt{degrees=5, perspective=0.0005}. \\
6 & Random erasing        & Partial occlusion by vegetation or shadow edges.
                           \texttt{erasing=0.3}. \\
7 & Copy-paste            & Places dummy instances onto new backgrounds.
                           \texttt{copy\_paste=0.1}. \\
8 & Fog/haze simulation   & UK weather; reduces contrast and saturation.
                           Grey-blending with $\alpha = 0.7$--$0.9$. \\
9 & Mosaic                & Tiles four images into one; critical for small-object
                           detection. \texttt{mosaic=1.0}. \\
10 & Scale variation       & Simulates altitude changes (15--50~m range).
                            \texttt{scale=0.7}. \\
\hline
\end{tabular}
\end{table}

Motion blur and brightness jitter together cover over 70\% of real-world failure
cases. Mosaic augmentation is particularly important for small-object detection,
as it forces the model to detect targets at various scales within a single training
image.

\subsection{Domain Gap Mitigation}

Several techniques reduce the gap between synthetic composites and real camera images:

\begin{itemize}
    \item \textbf{Composite edge blending:} Gaussian blur on the alpha mask prevents
          the model from learning ``sharp rectangular edge = target.''
    \item \textbf{Background-aware colour matching:} The dummy's mean brightness is
          matched to the local background patch where it is placed, preventing the
          model from exploiting brightness discontinuities.
    \item \textbf{Multiple background sources:} v1 used a single \texttt{map.jpg};
          v2 uses DJI video frames (dozens of distinct backgrounds). A minimum of
          20+ distinct backgrounds is recommended.
    \item \textbf{Real labelled frames:} The single most effective technique.
          Current v2 has 16 real frames; the target for v3 is 50--100.
\end{itemize}

\subsection{Dataset Composition Guidelines}

\begin{table}[htbp]
\centering
\caption{Recommended dataset composition for single-class aerial detection.}
\label{tab:dataset-composition}
\begin{tabular}{lll}
\hline
\textbf{Component} & \textbf{Ratio} & \textbf{Notes} \\
\hline
Synthetic images   & 50--75\% & Generated from \texttt{generate\_dataset\_v2.py} \\
Real labelled images & 25--50\% & From flight video via \texttt{label\_tool.py} \\
Negative images    & 15--25\% & Empty label files; reduces false positives \\
Hard negatives     & 5--10\% & Previous false positive frames \\
\hline
\end{tabular}
\end{table}


%% ======================================================================
\section{Lens Calibration}
\label{sec:lens-calibration}

\subsection{Checkerboard Calibration Procedure}

Lens distortion is calibrated using OpenCV's checkerboard corner detection with a
calib.io 14$\times$9 printed pattern (13$\times$8 inner corners, 28~mm square size).
The calibration script (\texttt{tests/calibration/lens\_calibrate.py}) supports both
interactive mode (with display) and headless mode (auto-capture every 3~seconds
for SSH operation).

The procedure captures 10--20 images of the checkerboard at various angles, distances,
and frame positions (including near-edge positions where distortion is worst). OpenCV
computes the camera intrinsic matrix $(f_x, f_y, c_x, c_y)$ and distortion coefficients
$(k_1, k_2, p_1, p_2, k_3)$ via \texttt{cv2.calibrateCamera()}.

\subsection{Undistortion Implementation}

When a valid \texttt{calibration\_data.npz} file exists and
\texttt{UNDISTORT\_ENABLED = True}, \texttt{VisionSystem} precomputes remap matrices
at startup using \texttt{cv2.initUndistortRectifyMap()} and applies
\texttt{cv2.remap()} to every frame before AI inference. This costs approximately
1.5~ms per frame on the Pi~5---negligible compared to the 206~ms inference time.

\subsection{When Undistortion Matters}

The IMX296 is a global shutter sensor with inherently low distortion. Bench testing
yielded an RMS reprojection error of 0.399. The practical impact is primarily on
detections near frame edges, where uncorrected barrel distortion shifts detection
positions by 1--3 pixels, translating to 1--2\% GPS estimation error at flight altitude.

For detections near frame centre (the most valuable observations due to the centrality
weighting scheme described in Section~\ref{sec:gps-estimation}), undistortion has
negligible effect. The system currently operates with \texttt{UNDISTORT\_ENABLED = False}
because the calibration file on the Pi was found to be corrupted during field testing;
the IMX296's minimal distortion makes this acceptable.


%% ======================================================================
\section{FOV Calibration}
\label{sec:fov-calibration}

\subsection{Focal Length Measurement}

The focal length is the single most important calibration parameter for GPS
estimation accuracy. A 10\% error in focal length produces a 10\% error in every
GPS estimate. Two calibration methods are available:

\textbf{Bench method (ruler):} The camera is placed at a known height $h$ pointing
straight down. The visible ground width $w_\text{vis}$ is measured with a tape measure.
The focal length is computed as:
\begin{equation}
    f_\text{mm} = \frac{w_\text{sensor} \times h}{w_\text{vis}}
    \label{eq:focal-length}
\end{equation}

\noindent where $w_\text{sensor} = 5.02$~mm is the IMX296 sensor width. At 1.0~m
height with 0.92~m visible width: $f = 5.02 \times 1.0 / 0.92 = 5.46$~mm (calibrated
2026-03-11).

\textbf{Video method (known object):} Using DJI flight video with SRT telemetry,
a known-size object (dummy, 1.8~m tall) is clicked at multiple altitudes. The
focal length in pixels is:
\begin{equation}
    f_\text{px} = \frac{\text{pixel\_height} \times \text{altitude}}
    {\text{known\_height}}
    \label{eq:focal-px}
\end{equation}

Nine samples across 15--50~m altitude yielded $f_\text{px} = 1416 \pm 41$~pixels,
corresponding to an HFOV of 54.4\textdegree\ for the 1456$\times$1088 DJI video crop.

\subsection{Ground Coverage Calculation}

The ground footprint at altitude $h$ is:
\begin{equation}
    W_\text{ground} = \frac{h \times w_\text{sensor}}{f_\text{mm}}
    = \frac{h \times 5.02}{5.46} \approx 0.92h
    \label{eq:ground-width}
\end{equation}

\begin{table}[htbp]
\centering
\caption{Ground coverage and ground sample distance (GSD) at operational altitudes.}
\label{tab:ground-coverage}
\begin{tabular}{rrrr}
\hline
\textbf{Altitude (m)} & \textbf{Ground width (m)} & \textbf{GSD (cm/px)} &
\textbf{Dummy height (px)} \\
\hline
10  & 9.2   & 0.63 & 286 \\
20  & 18.4  & 1.26 & 143 \\
25  & 23.0  & 1.58 & 114 \\
30  & 27.6  & 1.90 & 95 \\
35  & 32.2  & 2.21 & 81 \\
40  & 36.8  & 2.53 & 71 \\
50  & 46.0  & 3.16 & 57 \\
\hline
\end{tabular}
\end{table}

\subsection{Impact of Calibration Error}

Table~\ref{tab:fov-error} shows the GPS estimation error introduced by a 10\%
focal length miscalibration at various altitudes.

\begin{table}[htbp]
\centering
\caption{GPS position error from 10\% focal length miscalibration.}
\label{tab:fov-error}
\begin{tabular}{rr}
\hline
\textbf{Altitude (m)} & \textbf{Position error (m)} \\
\hline
10 & 0.9 \\
20 & 1.8 \\
30 & 2.7 \\
\hline
\end{tabular}
\end{table}


%% ======================================================================
\section{GPS Estimation from Camera Detections}
\label{sec:gps-estimation}

\subsection{Pixel-to-GPS Projection}

The core projection converts a detection's pixel coordinates into a GPS position
using the pinhole camera model. The pipeline proceeds in five steps:

\textbf{Step 1: Pixel offset from frame centre.}
\begin{align}
    \Delta x_\text{px} &= (x_\text{norm} - 0.5) \times W_\text{img} \\
    \Delta y_\text{px} &= (y_\text{norm} - 0.5) \times H_\text{img}
\end{align}

\noindent where $x_\text{norm} = c_x / W_\text{img}$ is the normalised detection
centre.

\textbf{Step 2: Convert pixels to metres on the ground.}
\begin{align}
    \Delta x_\text{m} &= \frac{\Delta x_\text{px} \times h}{f_\text{px}} \\
    \Delta y_\text{m} &= \frac{\Delta y_\text{px} \times h}{f_\text{px}}
\end{align}

\noindent where $h$ is altitude and $f_\text{px} = f_\text{mm} \times W_\text{img} / w_\text{sensor}$.

\textbf{Step 3: Map camera axes to world axes.}
\begin{align}
    d_\text{forward} &= -\Delta y_\text{m} \\
    d_\text{right} &= \Delta x_\text{m}
\end{align}

\textbf{Step 4: Rotate by drone yaw ($\psi$).}
\begin{align}
    d_\text{north} &= d_\text{forward} \cos\psi - d_\text{right} \sin\psi \\
    d_\text{east} &= d_\text{forward} \sin\psi + d_\text{right} \cos\psi
\end{align}

\textbf{Step 5: Convert to GPS offset.}
\begin{align}
    \hat{\phi} &= \phi_\text{drone} + \frac{d_\text{north}}{111{,}320} \\
    \hat{\lambda} &= \lambda_\text{drone} + \frac{d_\text{east}}
    {111{,}320 \cos \phi_\text{drone}}
\end{align}

\subsection{Inverse-Variance Weighting}

Each GPS observation receives a weight proportional to $1/h^2$, where $h$ is
altitude. This reflects the physics of perspective projection: ground sample
distance scales linearly with altitude, so estimation variance scales as $h^2$.
A detection at 10~m altitude receives 9$\times$ the weight of one at 30~m.

An additional centrality bonus of up to $5\times$ is applied for detections near
frame centre ($< 30$\% of frame diagonal from optical axis), where lens distortion
is minimal and the flat-earth projection is most accurate:
\begin{equation}
    w = \frac{1}{h^2} \times \left(1 + 4 \max\!\left(0,\;
    1 - \frac{d_\text{centre}}{0.3 \, d_\text{max}}\right)\right)
\end{equation}

The cumulative estimate is the weighted mean:
\begin{equation}
    \hat{\phi}_\text{target} = \frac{\sum_i w_i \hat{\phi}_i}{\sum_i w_i}
\end{equation}

In simulation testing, inverse-variance weighting achieved approximately 0.5~m
estimation error after 30~seconds of hovering at 10--15~m, compared to approximately
2~m with equal weighting.

\subsection{SmartEstimator: Tightest-Cluster Lock}

The \texttt{DummyEstimator} weighted mean is susceptible to outliers from false
positives or GPS timing lag. The \texttt{SmartEstimator} addresses this by finding
the tightest cluster of $N$ consistent estimates using a greedy medoid-expansion
algorithm:

\begin{enumerate}
    \item Compute all pairwise GPS distances between accumulated estimates.
    \item Identify the geometric medoid (point with smallest total distance to all others).
    \item Greedily add the nearest remaining point, recomputing cluster spread.
    \item Stop at $N$ points (default $N = 5$, configurable via \texttt{-{}-smart-min}).
    \item If the maximum pairwise distance within the cluster is less than
          \texttt{max\_spread} (default 1.0~m), the estimator \textbf{locks}.
\end{enumerate}

Once locked, no further estimates are accepted. The final coordinate is the
\textbf{median} latitude and longitude of the locked cluster (robust to within-cluster
outliers). The algorithm is deterministic and $O(N^2)$, which is negligible for
typical mission sizes of 10--50 detections per target.

\subsection{GPS Timing Lag}

The GPS receiver has 100--200~ms latency. At 5~m/s flight speed, this produces
0.5--1.0~m systematic position error along the flight direction. Analysis of DJI
flight video (CEP50 = 2.3~m, maximum error 16.5~m) confirmed this effect.
Multiple passes from different headings average out the directional bias. The
SmartEstimator's cluster median further mitigates the effect by selecting the
most consistent subset of estimates.


%% ======================================================================
\section{Inference Performance}
\label{sec:inference-performance}

\subsection{Raspberry Pi 5 Benchmarks}

All benchmarks were conducted on the Raspberry Pi~5 (Broadcom BCM2712, 4-core
Cortex-A76 at 2.4~GHz) using the TFLite XNNPACK CPU delegate. Results are
summarised in Table~\ref{tab:pi-benchmarks}.

\begin{table}[htbp]
\centering
\caption{Inference benchmarks on Raspberry Pi~5.}
\label{tab:pi-benchmarks}
\begin{tabular}{lrrrr}
\hline
\textbf{Backend} & \textbf{Avg time (ms)} & \textbf{FPS} &
\textbf{Detection rate} & \textbf{Confidence} \\
\hline
TFLite (float32)          & 206.5 & 4.8  & 50/50 & 0.966 \\
TFLite + undistortion     & 208.0 & 4.8  & 50/50 & 0.958 \\
NCNN (expected)           & $\sim$72 & $\sim$13.8 & --- & --- \\
\hline
\end{tabular}
\end{table}

The undistortion overhead is approximately 1.5~ms per frame, negligible relative
to inference time. All models have identical input/output shapes and are drop-in
replacements.

\subsection{Speed Improvement Options}

Research into inference acceleration identified the following ranked options:

\begin{table}[htbp]
\centering
\caption{CV inference speed improvement options for Pi~5.}
\label{tab:speed-options}
\begin{tabular}{clrl}
\hline
\textbf{Rank} & \textbf{Method} & \textbf{Expected FPS} & \textbf{Effort} \\
\hline
1 & FP16 XNNPACK          & $\sim$10 & Low \\
2 & Threaded pipeline      & +20--30\% & Medium \\
3 & Lower input resolution & 8--20    & Low \\
4 & NCNN backend           & $\sim$15 & Medium \\
5 & Overclock to 2.8~GHz   & +15\%    & Trivial \\
\hline
\end{tabular}
\end{table}

\subsection{Frames per Target Pass}

At the search altitude of 35~m, the ground footprint is approximately 32~m wide.
At 10~m/s flight speed, the drone traverses the footprint in 3.2~seconds, yielding
approximately 15 frames per target pass at 4.8~FPS. The required per-frame
detection probability for 95\% pass detection is:
\begin{equation}
    p > 1 - 0.05^{1/N}
    \label{eq:detection-probability}
\end{equation}

For $N = 15$ frames, $p > 0.18$ suffices for 95\% detection probability across
the pass. The current model achieves effectively 100\% per-frame detection on
the test image (confidence 0.966), providing substantial margin.


%% ======================================================================
\section{Detection Reliability}
\label{sec:detection-reliability}

\subsection{Detection Altitude Ceiling}

Expected target size in the YOLO model input space (640$\times$640) determines
the detection ceiling. Table~\ref{tab:altitude-detection} summarises the relationship.

\begin{table}[htbp]
\centering
\caption{Dummy size in model input space and expected detection performance.}
\label{tab:altitude-detection}
\begin{tabular}{rrrp{4cm}}
\hline
\textbf{Altitude (m)} & \textbf{Frame pixels} & \textbf{Model pixels} &
\textbf{Detection likelihood} \\
\hline
20 & $\sim$67 & $\sim$29 & Very high ($>$99\%) \\
25 & $\sim$54 & $\sim$24 & Very high ($>$98\%) \\
30 & $\sim$45 & $\sim$20 & High ($>$95\%) \\
35 & $\sim$38 & $\sim$17 & High ($>$90\%) \\
40 & $\sim$33 & $\sim$15 & Moderate (70--90\%) \\
50 & $\sim$27 & $\sim$12 & Uncertain (50--80\%) \\
\hline
\end{tabular}
\end{table}

Below approximately 10 pixels in model input space, YOLO detection degrades rapidly.
The practical ceiling is estimated at 40--45~m for reliable detection ($>$90\%
per-frame). At the configured search altitude of 35~m, the dummy occupies
approximately 17 pixels in model space---above the critical threshold but without
large margin.

\subsection{False Positive Mitigation}

False positives arise from shadows, rocks, fence posts, tractor tracks, and field
boundary patterns. The system employs a four-layer defence:

\begin{enumerate}
    \item \textbf{YOLO inference} with low threshold (0.2): catch all potential
          targets, accept false positives at this stage.
    \item \textbf{Post-detection filters}: size filtering (reject bounding boxes
          inconsistent with a 1.8~m dummy at current altitude, expected 30--50\%
          FP reduction), aspect ratio filtering (reject near-square or extremely
          elongated detections, 15--25\% FP reduction), and edge proximity rejection
          (reject detections within 5\% of frame edges, 10--15\% FP reduction).
    \item \textbf{Temporal voting}: require 2 out of 5 consecutive frames to
          produce detections at a consistent GPS location (within 5~m radius).
          Expected 40--60\% FP reduction. Transient false positives from
          single-frame noise are completely eliminated.
    \item \textbf{Operator confirmation}: Y/N/I/X buttons in the ground station
          browser interface. The human safety net.
\end{enumerate}

With all automated layers active, the expected false positive rate is less than
1 per complete scan. The true positive rate target is $>$95\% for targets within
the detection altitude ceiling.

\subsection{Hard Negative Mining}

Each flight improves the model through hard negative mining:
\begin{enumerate}
    \item The operator presses X (false positive) during flight when the system
          triggers on a non-target.
    \item The raw frame is saved automatically as a training negative.
    \item After the flight, these frames are added to the training dataset with
          empty label files and the model is retrained.
\end{enumerate}

Expected impact: 30--50\% false positive reduction per retraining cycle. After
2--3 cycles, the model learns site-specific visual confusers.

\subsection{Confidence Threshold Strategy}

The confidence threshold of 0.2 is deliberately low. The reasoning: missing the
actual target is catastrophic, while extra false positives are merely inconvenient
(filtered by the operator). At maximum altitude (50~m), the dummy occupies only
$\sim$12 pixels in model space, and confidence may drop to 0.3--0.5. Raising
the threshold above 0.4 risks missing detections at the altitude ceiling.

A two-threshold strategy is planned for future implementation:
\begin{itemize}
    \item Search pass (35~m): threshold 0.15 (wide net, log all candidates).
    \item Rescan pass (28~m, 22~m): threshold 0.30 (revisit only promising locations).
    \item Verify hover (15~m): threshold 0.40 (close-up confirmation).
\end{itemize}


%% ======================================================================
\section{Model Deployment}
\label{sec:model-deployment}

\subsection{TFLite Export}

The primary deployment format is TFLite float32, exported from the trained PyTorch
weights:

\begin{verbatim}
model = YOLO('runs/train/sar_v2/weights/best.pt')
model.export(format='tflite', imgsz=640)
\end{verbatim}

\noindent The export always uses \texttt{imgsz=640} regardless of training resolution,
because \texttt{vision.py} resizes all frames to 640$\times$640 before inference.
The resulting file (\texttt{best\_float32.tflite}, $\sim$11.7~MB) is a complete
self-contained inference graph including NMS.

\subsection{NCNN Export}

For approximately $3\times$ faster inference on the Pi~5's ARM Cortex-A76 cores:

\begin{verbatim}
model.export(format='ncnn', imgsz=640, half=True)
\end{verbatim}

\noindent The NCNN export produces a directory containing \texttt{model.ncnn.bin}
and \texttt{model.ncnn.param} files ($\sim$7~MB total). The FP16 (half-precision)
option leverages the Pi~5's native FP16 hardware support. Vision.py loads NCNN
models via the \texttt{ncnn} Python package with 4-thread inference and a warmup
pass to trigger JIT compilation.

\subsection{Model Swapping Workflow}

All scripts read \texttt{best.tflite} from the project root. Swapping models
requires no code changes:

\begin{verbatim}
# Option A: Copy model to active slot
cp cv_models/sar_v2_1088/best.tflite best.tflite

# Option B: Specify at runtime
python main.py --model cv_models/sar_v2_1088/best.tflite
\end{verbatim}

Any TFLite model with input shape $[1, 640, 640, 3]$ and output shape
$[1, 5, 8400]$ is a drop-in replacement. Available model variants include:

\begin{table}[htbp]
\centering
\caption{Available model variants for deployment.}
\label{tab:model-variants}
\begin{tabular}{llrl}
\hline
\textbf{Model} & \textbf{Location} & \textbf{Size} & \textbf{Notes} \\
\hline
Custom dummy (v2) & \texttt{cv\_models/sar\_v2\_1088/} & 11.7~MB & Best, mAP50=0.995 \\
Custom dummy (v1) & \texttt{models/custom\_yolov8n.tflite} & 3.3~MB & Baseline \\
COCO person (80-class) & \texttt{models/human.tflite} & 13~MB & Backup/fallback \\
NCNN (v2)         & \texttt{cv\_models/sar\_v2\_1088/ncnn/} & 7~MB & Fastest on Pi \\
\hline
\end{tabular}
\end{table}

\subsection{Deployment Pipeline}

The complete deployment workflow from training to field operation:

\begin{enumerate}
    \item Train on Google Colab (T4 GPU, 30--60 minutes).
    \item Export TFLite (\texttt{imgsz=640}) and optionally NCNN.
    \item Download and save to \texttt{cv\_models/sar\_vN/}.
    \item Test on laptop with DJI video replay
          (\texttt{tests/laptop/video\_test.py}).
    \item Benchmark on Pi (\texttt{tests/hardware/benchmark.py}): verify
          $\sim$200~ms inference, 50/50 detection, $>$0.9 confidence.
    \item Deploy: \texttt{cp cv\_models/sar\_vN/best.tflite best.tflite},
          push to GitHub, pull on Pi.
    \item Validate with passive watch during manual flight.
\end{enumerate}

This pipeline ensures that every model is tested against real flight footage on the
laptop and benchmarked on the target hardware before field deployment. The
\texttt{-{}-model} CLI flag enables rapid A/B comparison without file copying.


%% ======================================================================
\section{Preprocessing Pipeline}
\label{sec:preprocessing}

Table~\ref{tab:preprocessing} summarises the available preprocessing techniques
and their status.

\begin{table}[htbp]
\centering
\caption{Preprocessing techniques and their deployment status.}
\label{tab:preprocessing}
\begin{tabular}{lrlp{4.5cm}}
\hline
\textbf{Technique} & \textbf{Cost (ms)} & \textbf{Default} & \textbf{When to enable} \\
\hline
Lens undistortion & 1.5 & OFF & When valid calibration file exists \\
CLAHE             & 2--3 & OFF & Long shadows, mixed lighting \\
Grey-world AWB    & $<$1 & OFF & Persistent colour cast \\
Sharpening        & $<$1 & OFF & Only if measurably helps detection \\
Contrast stretching & 1--2 & OFF & Haze, fog, mist \\
\hline
\end{tabular}
\end{table}

The total worst-case overhead with all preprocessing enabled is approximately 6~ms,
representing only 3\% of the 206~ms inference time---negligible impact on the
detection pipeline.

CLAHE (Contrast Limited Adaptive Histogram Equalisation) is the most promising
optional preprocessing step for the UK operational environment. It normalises
local contrast across the frame, making targets visible in both shadowed and
sun-bleached regions simultaneously. Applied in the LAB colour space (L channel
only), it adds 2--3~ms per frame and can improve detection confidence by
5--15\% in mixed-lighting conditions.
